% Options for packages loaded elsewhere
\PassOptionsToPackage{unicode}{hyperref}
\PassOptionsToPackage{hyphens}{url}
\PassOptionsToPackage{dvipsnames,svgnames,x11names}{xcolor}
\documentclass[
  11pt,
]{article}
\usepackage{xcolor}
\usepackage[margin=1in]{geometry}
\usepackage{amsmath,amssymb}
\setcounter{secnumdepth}{-\maxdimen} % remove section numbering
\usepackage{iftex}
\ifPDFTeX
  \usepackage[T1]{fontenc}
  \usepackage[utf8]{inputenc}
  \usepackage{textcomp} % provide euro and other symbols
\else % if luatex or xetex
  \usepackage{unicode-math} % this also loads fontspec
  \defaultfontfeatures{Scale=MatchLowercase}
  \defaultfontfeatures[\rmfamily]{Ligatures=TeX,Scale=1}
\fi
\usepackage{lmodern}
\ifPDFTeX\else
  % xetex/luatex font selection
  \setmainfont[]{Libertinus Serif}
  \setmathfont[]{Libertinus Math}
\fi
% Use upquote if available, for straight quotes in verbatim environments
\IfFileExists{upquote.sty}{\usepackage{upquote}}{}
\IfFileExists{microtype.sty}{% use microtype if available
  \usepackage[]{microtype}
  \UseMicrotypeSet[protrusion]{basicmath} % disable protrusion for tt fonts
}{}
\makeatletter
\@ifundefined{KOMAClassName}{% if non-KOMA class
  \IfFileExists{parskip.sty}{%
    \usepackage{parskip}
  }{% else
    \setlength{\parindent}{0pt}
    \setlength{\parskip}{6pt plus 2pt minus 1pt}}
}{% if KOMA class
  \KOMAoptions{parskip=half}}
\makeatother
\usepackage{longtable,booktabs,array}
\newcounter{none} % for unnumbered tables
\usepackage{calc} % for calculating minipage widths
% Correct order of tables after \paragraph or \subparagraph
\usepackage{etoolbox}
\makeatletter
\patchcmd\longtable{\par}{\if@noskipsec\mbox{}\fi\par}{}{}
\makeatother
% Allow footnotes in longtable head/foot
\IfFileExists{footnotehyper.sty}{\usepackage{footnotehyper}}{\usepackage{footnote}}
\makesavenoteenv{longtable}
\setlength{\emergencystretch}{3em} % prevent overfull lines
\providecommand{\tightlist}{%
  \setlength{\itemsep}{0pt}\setlength{\parskip}{0pt}}
\usepackage{longtable,booktabs,array}
\usepackage{ragged2e}
\usepackage{makecell}
\AtBeginEnvironment{longtable}{\small}

\usepackage{amsmath,amssymb,mathtools}
\usepackage{microtype}
\usepackage[none]{hyphenat}
\emergencystretch=3em

% Centered publication title block (Pandoc \maketitle)
\usepackage{titling}
\pretitle{\begin{center}\LARGE\bfseries}
\posttitle{\par\end{center}\vskip 0.5em}
\preauthor{\begin{center}\large}
\postauthor{\par\end{center}}
\predate{\begin{center}\large}
\postdate{\par\end{center}\vskip 1.5em}

% Allow breaks in pandoc inline-code paths (Appendix B filenames)
\usepackage{xurl}
\makeatletter
\def\UrlBreaks{%
  \do\-\do\.\do\_\do\/\do\=\do\?\do\&\do\ \do\a\do\b\do\c\do\d\do\e\do\f\do\g\do\h\do\i\do\j%
  \do\k\do\l\do\m\do\n\do\o\do\p\do\q\do\r\do\s\do\t\do\u\do\v\do\w\do\x\do\y\do\z%
  \do\A\do\B\do\C\do\D\do\E\do\F\do\G\do\H\do\I\do\J\do\K\do\L\do\M\do\N\do\O\do\P\do\Q\do\R\do\S\do\T\do\U\do\V\do\W\do\X\do\Y\do\Z\do\0\do\1\do\2\do\3\do\4\do\5\do\6\do\7\do\8\do\9}
\makeatother
\usepackage{bookmark}
\IfFileExists{xurl.sty}{\usepackage{xurl}}{} % add URL line breaks if available
\urlstyle{same}
\hypersetup{
  colorlinks=true,
  linkcolor={blue},
  filecolor={Maroon},
  citecolor={Blue},
  urlcolor={blue},
  pdfcreator={LaTeX via pandoc}}

\author{}
\date{}

\begin{document}

\section{The Shape of Time}\label{the-shape-of-time}

\section{Appendix A --- Popular Paper, Reader Edition (NTE v10.8
v7)}\label{appendix-a-popular-paper-reader-edition-nte-v10.8-v7}

\emph{Non-Temporal Emergence (NTE): the Null-Time Extension to the
Standard Model}

G. White

\begin{center}\rule{0.5\linewidth}{0.5pt}\end{center}

This appendix is \textbf{not} the full corpus Popular Paper. It
reproduces the \textbf{Reader Edition} only:

{\def\LTcaptype{none} % do not increment counter
\begin{longtable}[]{@{}
  >{\raggedright\arraybackslash}p{(\linewidth - 4\tabcolsep) * \real{0.4286}}
  >{\raggedright\arraybackslash}p{(\linewidth - 4\tabcolsep) * \real{0.2857}}
  >{\raggedright\arraybackslash}p{(\linewidth - 4\tabcolsep) * \real{0.2857}}@{}}
\toprule\noalign{}
\begin{minipage}[b]{\linewidth}\raggedright
Edition
\end{minipage} & \begin{minipage}[b]{\linewidth}\raggedright
File
\end{minipage} & \begin{minipage}[b]{\linewidth}\raggedright
Role
\end{minipage} \\
\midrule\noalign{}
\endhead
\bottomrule\noalign{}
\endlastfoot
\textbf{Reader Edition (this appendix)} &
\texttt{Paper\_\textbackslash{}allowbreak\{\}Popular\_\textbackslash{}allowbreak\{\}v10\_\textbackslash{}allowbreak\{\}8\_\textbackslash{}allowbreak\{\}v7.md}
& \textasciitilde4,800 words; recap after the chapters \\
\textbf{Full Edition (corpus)} &
\texttt{Paper\_\textbackslash{}allowbreak\{\}Popular\_\textbackslash{}allowbreak\{\}v10\_\textbackslash{}allowbreak\{\}8\_\textbackslash{}allowbreak\{\}v6.md}
& \textasciitilde13,000 words; complete corpus lay paper \\
\end{longtable}
}

Same framework and predictions; different length and purpose. Technical
proofs and companion papers: \textbf{Appendix B}.

\textbf{Source of truth for this appendix:}
\texttt{NTE\ v10\_\textbackslash{}allowbreak\{\}6/Papers/Paper\_\textbackslash{}allowbreak\{\}Popular\_\textbackslash{}allowbreak\{\}v10\_\textbackslash{}allowbreak\{\}8\_\textbackslash{}allowbreak\{\}v7.md}

\begin{center}\rule{0.5\linewidth}{0.5pt}\end{center}

\section{What If Time Is Made of
Something?}\label{what-if-time-is-made-of-something}

\subsection{The Null-Time Extension to the Standard Model, Explained
From First
Principles}\label{the-null-time-extension-to-the-standard-model-explained-from-first-principles}

\textbf{G. White}

\begin{center}\rule{0.5\linewidth}{0.5pt}\end{center}

\begin{quote}
\textbf{NTE Popular Paper --- v10.8 v7 · Reader Edition (book
appendix)}\\
\textasciitilde4,800 words. Compressed lay recap for readers of
\emph{The Shape of Time}; canonical source for \textbf{Appendix A}
(\texttt{Appendix\_\textbackslash{}allowbreak\{\}A\_\textbackslash{}allowbreak\{\}NTE\_\textbackslash{}allowbreak\{\}Popular\_\textbackslash{}allowbreak\{\}Paper\_\textbackslash{}allowbreak\{\}v10\_\textbackslash{}allowbreak\{\}8\_\textbackslash{}allowbreak\{\}v7.md}).\\
\textbf{This is not the full corpus paper.} The companion document is
\textbf{\texttt{Paper\_\textbackslash{}allowbreak\{\}Popular\_\textbackslash{}allowbreak\{\}v10\_\textbackslash{}allowbreak\{\}8\_\textbackslash{}allowbreak\{\}v6.md}
--- Full Edition} (\textasciitilde13,000 words): the complete corpus lay
paper for standalone distribution. Same framework; different length and
purpose. Do not treat v6 and v7 as interchangeable filenames.
\end{quote}

\begin{center}\rule{0.5\linewidth}{0.5pt}\end{center}

There is a fact buried in your physics textbook that nobody has fully
followed to its conclusion.

A photon of light --- any photon, ever --- experiences zero time between
the moment it's created and the moment it's absorbed. Not very little
time. Zero. From the photon's perspective (to the extent that word means
anything for an object with no clock), emission and absorption are the
same instant. The entire journey from a star four hundred light-years
away to your eye happens, for the photon, in no time at all.

Every physicist knows this. It follows directly from Einstein's special
relativity, one of the most tested theories in all of science. And then,
almost universally, physicists set it aside and get on with their
calculations.

This paper is what happens when you don't set it aside.

\begin{center}\rule{0.5\linewidth}{0.5pt}\end{center}

\subsection{A Note on Method}\label{a-note-on-method}

This paper is not a claim that NTE is the correct theory of fundamental
physics. It's a working bet --- the kind of bet Feynman described in
\emph{QED} when he asked why all theories of physics look so similar in
structure. He named three possibilities: that physicists are limited in
imagination; that Nature has only one way of doing things and repeats
her story; or that apparent similarities are fingers of the same hand.

NTE bets on the second and third --- that when multiple mathematical
descriptions agree on a phenomenon, the agreement is evidence of common
structure, and that unification means finding the structure rather than
choosing between descriptions.

This is a bet, not a law. It can be wrong. Success or failure will be
measured the way any framework is measured: predictions that hold,
derivations that survive scrutiny, and whether the structure keeps
simplifying as observations arrive. Beauty is a guide to which path to
try first, not a shortcut past empirical test.

The rest of this paper is what the bet has produced so far.

\begin{center}\rule{0.5\linewidth}{0.5pt}\end{center}

\subsection{The Clue Nobody Followed
Up}\label{the-clue-nobody-followed-up}

Photons are massless. Every massless entity in the universe experiences
exactly zero proper time. Every massive entity accumulates proper time.
The correlation is exact: zero exceptions.

The Standard Model explains mass through the Higgs field at CERN. NTE
reads the mass--time correlation as a \emph{surface signature} of
something deeper: two coordinated roles set at the same cosmic moment,
not one mechanism doing both jobs.

{\def\LTcaptype{none} % do not increment counter
\begin{longtable}[]{@{}
  >{\raggedright\arraybackslash}p{(\linewidth - 4\tabcolsep) * \real{0.2500}}
  >{\raggedright\arraybackslash}p{(\linewidth - 4\tabcolsep) * \real{0.3056}}
  >{\raggedright\arraybackslash}p{(\linewidth - 4\tabcolsep) * \real{0.4444}}@{}}
\toprule\noalign{}
\begin{minipage}[b]{\linewidth}\raggedright
What you observe
\end{minipage} & \begin{minipage}[b]{\linewidth}\raggedright
Standard Model Higgs
\end{minipage} & \begin{minipage}[b]{\linewidth}\raggedright
Boundary expansion field (NTE)
\end{minipage} \\
\midrule\noalign{}
\endhead
\bottomrule\noalign{}
\endlastfoot
Mass & Yukawa couplings to fermions & Not its job \\
Proper time & Not its job & How fast expansion becomes experienced
time \\
Photon & Zero Higgs coupling & Zero engagement with expansion field \\
Heavy quark & Large Yukawa & Deep engagement with expansion field \\
\end{longtable}
}

\textbf{Non-Temporal Emergence (NTE)} --- the \textbf{Null-Time
Extension to the Standard Model} --- adds one geometric scalar, written
\textbar Φ\_NTE\textbar, on the universe's expanding boundary. It does
not replace the Standard Model. The SM Higgs still gives mass;
\textbar Φ\_NTE\textbar{} governs how cosmic expansion accumulates
proper time for matter on the surface. Both identities were assigned in
one coordinated cascade at \emph{gauge birth}: the moment gauge symmetry
and clocks switched on. The photon landed at the zero-coupling place on
both sides.

\begin{center}\rule{0.5\linewidth}{0.5pt}\end{center}

\subsection{The Universe as an Expanding
Surface}\label{the-universe-as-an-expanding-surface}

Forget a universe exploding into pre-existing void. In NTE, our
three-dimensional cosmos is the \emph{surface} of a four-dimensional
sphere --- like an ant on an inflating balloon who cannot leave the
rubber. Unlike a balloon, there is no ``inside'' to visit: only the
surface and its expansion history exist as places.

Expansion is not measured first in seconds. It's measured in a
dimensionless count \emph{N} --- how many times the universe has doubled
in size. That count plays the role time plays in standard cosmology,
without presupposing clocks.

The local expansion rate ties to the boundary field: \textbf{H =
α\textbar Φ\_NTE\textbar²}. Deeper engagement with
\textbar Φ\_NTE\textbar{} means deeper embedding in time ---
gravitational time dilation in everyday language.

\textbf{Why are we not moving at the speed of light?} Because photons
ride the expansion wavefront at \emph{c} with zero field engagement,
while massive matter is embedded slightly behind the wavefront. \emph{c}
is the manifold's expansion rate per unit accumulated time, not merely a
photon property. Distant galaxies beyond the Hubble radius recede faster
than \emph{c} without violating relativity: that is geometry on an
expanding surface, not objects outrunning light through space.

\begin{center}\rule{0.5\linewidth}{0.5pt}\end{center}

\subsection{One Manifold, Two
Orientations}\label{one-manifold-two-orientations}

At the Janus boundary, the universe acquired bilateral structure: one
manifold with \emph{thickness} along the e-fold direction, not two
separate bubbles. We occupy one orientation; the conjugate orientation
is CPN-reversed (matter, parity, and e-fold direction flipped). Neither
side is privileged. Gravitational time dilation is how that thickness
varies where \textbar Φ\_NTE\textbar{} is weaker in potential wells.

\begin{center}\rule{0.5\linewidth}{0.5pt}\end{center}

\subsection{Entanglement: the Photon's Trick on a Larger
Scale}\label{entanglement-the-photons-trick-on-a-larger-scale}

If emission and absorption can be the same event for a photon on the
wavefront, the same geometry may explain entanglement.

Two entangled particles are not two objects needing a hidden signal.
They are endpoints of one zero-proper-time structure on the surface ---
like a photon's path, but with internal structure. Correlations are
properties of the shape, not messages sent between separated things.

Bell's bound (2) and quantum mechanics' Tsirelson ceiling (2√2) are
often treated as unexplained constants. Here 2√2 is read as angular
overlap on the three-sphere through the Hopf fibration --- consistent
with saturating experiments, not proof by itself.

\begin{center}\rule{0.5\linewidth}{0.5pt}\end{center}

\subsection{The Neutrino That Bridges Both
Sides}\label{the-neutrino-that-bridges-both-sides}

The neutrino is the Standard Model's odd particle: tiny mass, feeble
interactions, passes through Earth untouched.

In NTE, every other massive species completed \emph{Hopf-projection} at
gauge birth --- a topological lock to one orientation. The neutrino's
lock did not finish because the Majorana particle--antiparticle binary
is still unresolved topologically (not because the Yukawa coupling is
``too small''). What survives is a chirality superposition spanning both
orientations of the manifold's thickness.

Why do weak experiments see only left-handed neutrinos? Each orientation
samples the same underlying rhythm at opposite phase; our weak force
catches the left-handed phase. Parity violation in the weak sector
becomes a sampling asymmetry, not a brute fact.

That bilateral support makes the neutrino the backbone for
cross-thickness gravity.

\begin{center}\rule{0.5\linewidth}{0.5pt}\end{center}

\subsection{Dark Matter as Cross-Thickness
Gravity}\label{dark-matter-as-cross-thickness-gravity}

Matter in the conjugate orientation is invisible electromagnetically but
should gravitate into ours through the neutrino bridge. \textbf{Dark
matter is that gravitational signature}, not a new particle hiding in
underground tanks.

\textbf{κ ≈ 5.50} --- the dark-to-ordinary matter ratio --- comes from
one structural derivation in two steps: a bare coupling κ₀ = 3/√(4π)
from the unresolved neutrino's geometry on the three-sphere, then a
self-consistent resummation κ = κ₀/(1 − κ₀). Planck measures 5.38 ± 0.15
(within \textasciitilde0.8σ). Direct-detection null results are
predicted, not embarrassing failures. Dwarf galaxies should show scatter
in dark-matter fractions; large clusters should converge to 5.50 --- set
by neutrino free-streaming coherence (\textasciitilde40 Mpc).

\begin{center}\rule{0.5\linewidth}{0.5pt}\end{center}

\subsection{The Arrow of Time}\label{the-arrow-of-time}

Fundamental equations are time-symmetric; experience is not. Boltzmann's
entropy story explains direction only after assuming an extraordinarily
special initial state.

NTE locates the arrow in geometry: at \textbf{N\_lock}, when gauge
symmetry switched on, massive matter committed irreversibly to one
orientation of the manifold's thickness. The past is expanded e-fold
structure; the future is what has not yet expanded. Entropy still rises,
but the deeper asymmetry is geometric commitment at birth --- closer to
Pratchett's \emph{Thief of Time} (\emph{everything that happens, stays
happened}) than to statistics alone.

\begin{center}\rule{0.5\linewidth}{0.5pt}\end{center}

\subsection{Why Quantum Mechanics Is
Geometry}\label{why-quantum-mechanics-is-geometry}

Quantum states live on a three-sphere with unobservable global phase.
Quotienting phase via the Hopf fibration leaves probabilities on a
two-sphere. The Born rule \(P = \left|\langle\psi|\phi\rangle\right|^2\)
is then the unique smooth, symmetric, normalised choice on that geometry
--- not an extra postulate smuggled in.

The exponent is 2 because the defining constraint on complex amplitudes
is quadratic. Recent technical work sharpens this once gauge birth
activates the U(1) fibre; deriving complex structure from pre-birth real
data without assuming \(\mathbb{C}\) upfront remains open work.

\begin{center}\rule{0.5\linewidth}{0.5pt}\end{center}

\subsection{The Moment Physics Came Into
Being}\label{the-moment-physics-came-into-being}

Before \textbf{N\_lock}, there was no well-defined dt/dN and no Standard
Model Higgs as we know it --- no photons, no atoms, no us. The dominant
pre-birth dynamics live in a \textbf{bulk} geometry; our universe
\textbf{crystallises on the boundary} at gauge birth when the constraint
H = α\textbar Φ\_NTE\textbar² activates. That transition releases
\textasciitilde10³² K --- the Big Bang as physics switching on, not an
explosion into pre-existing space.

Pre-birth content: topological spinor modes on S³ --- seeds without
mass, charge, or interactions.

\begin{center}\rule{0.5\linewidth}{0.5pt}\end{center}

\subsection{The Algebra of Creation}\label{the-algebra-of-creation}

Pure mathematics (Hurwitz) gives exactly four division algebras usable
for quantum amplitudes: \(\mathbb{R}\), \(\mathbb{C}\), \(\mathbb{H}\),
\(\mathbb{O}\). Pre-temporal physics is octonionic (non-associative ---
no ordering of operations without time). When time orders operations,
associativity forces a drop to quaternions; the residual symmetry is
\textbf{SU(3) × SU(2) × U(1)} --- the Standard Model gauge group. The
detailed steps live in the technical corpus; the popular claim is
structural: the SM group is what time \emph{forces}, not an arbitrary
dial.

\begin{center}\rule{0.5\linewidth}{0.5pt}\end{center}

\subsection{Hubble Tension and Dark Energy as One
Phenomenon}\label{hubble-tension-and-dark-energy-as-one-phenomenon}

Local supernova surveys (\textasciitilde73 km/s/Mpc) disagree with CMB
back-calculation (\textasciitilde67). DESI hints dark energy's equation
of state \emph{w} may sit slightly above −1 at low redshift.

NTE reads both as \textbf{void bias}: most cosmic volume is empty;
\textbar Φ\_NTE\textbar{} is higher in voids, so lines of sight through
voids see slightly faster local expansion. CMB integrates the cosmic
mean; local candles prefer voids. True \emph{w} = −1 exactly; apparent
evolution is a measurement bias from spatial variation of
\textbar Φ\_NTE\textbar², not new dynamical dark energy.

\begin{center}\rule{0.5\linewidth}{0.5pt}\end{center}

\subsection{Prediction Scorecard}\label{prediction-scorecard}

{\def\LTcaptype{none} % do not increment counter
\begin{longtable}[]{@{}
  >{\raggedright\arraybackslash}p{(\linewidth - 6\tabcolsep) * \real{0.1892}}
  >{\raggedright\arraybackslash}p{(\linewidth - 6\tabcolsep) * \real{0.1351}}
  >{\raggedright\arraybackslash}p{(\linewidth - 6\tabcolsep) * \real{0.5135}}
  >{\raggedright\arraybackslash}p{(\linewidth - 6\tabcolsep) * \real{0.1622}}@{}}
\toprule\noalign{}
\begin{minipage}[b]{\linewidth}\raggedright
Claim
\end{minipage} & \begin{minipage}[b]{\linewidth}\raggedright
NTE
\end{minipage} & \begin{minipage}[b]{\linewidth}\raggedright
Observation / test
\end{minipage} & \begin{minipage}[b]{\linewidth}\raggedright
Tier
\end{minipage} \\
\midrule\noalign{}
\endhead
\bottomrule\noalign{}
\endlastfoot
Dark matter ratio κ & 5.50 & 5.38 ± 0.15 (Planck) & Clean
zero-parameter \\
Direct DM detection & Null forever & XENON, LUX, PandaX nulls &
Falsifiable \\
Neutrino nature & Majorana (no completed Hopf lock) & Neutrinoless
double beta decay & Falsifiable \\
Neutrino generations & Exactly 3 & No sterile fourth & Structural \\
m\_\{ν₂\} & \textasciitilde0.00860 eV & √(Δm²\_solar)
\textasciitilde0.00866 eV & Clean zero-parameter \\
n\_s & 0.9685 & 0.9649 ± 0.0042 & Clean zero-parameter \\
r (hilltop) & \textasciitilde4×10⁻³ & CMB-S4 / LiteBIRD & Future \\
Λ & \textasciitilde1.05× observed & Precision cosmology & Real test (α
anchoring caveat) \\
H₀ cosmic mean & 67.6 km/s/Mpc & CMB 67.4; local \textasciitilde73
void-biased & Structural \\
w & −1 exactly & DESI apparent w \textgreater{} −1 at low z & Structural
bias \\
θ₂₃ & 45° & 45.0° ± 1.0° & Clean \\
sin²θ₁₂ & 1/3 leading & \textasciitilde0.307 & Partial \\
Baryon asymmetry & Exists; matter sign fixed & η \textasciitilde{}
6.1×10⁻¹⁰ magnitude open & Mixed \\
W\_R boson & None & LHC nulls continue & Distinguishing \\
Quark/lepton mass ratios & Not derived & Inherited like SM & Open \\
\end{longtable}
}

\textbf{Six-from-one-chord:} One heat-kernel reading on the
universe-instrument's spectral triple lands six slots: κ₀
-\textgreater{} κ, gauge-birth temperature, inflation tilt (n\_s and r
on the hilltop reading), Majorana and leptogenesis scales, Λ ≈ 1.05×
observed, and the comoving-volume audit. Chapter 12 develops the chord;
Appendix B carries the technical derivations.

For the Λ note in lay terms: the framework first removes the
catastrophic vacuum term by the volume constraint, then keeps the small
spectral leftover, then renormalises that leftover over the comoving
volume at gauge birth using the same long e-fold depth as n\_s (not the
short lock width). That is where the \textasciitilde1.05× value comes
from.

\begin{center}\rule{0.5\linewidth}{0.5pt}\end{center}

\subsection{Two Registers Before Physics Switched
On}\label{two-registers-before-physics-switched-on}

The CMB tilt \textbf{n\_s ≈ 0.9685} comes from a long pre-birth count
(\textasciitilde63 e-folds) in the \emph{ripple register}. The lock to
gauge birth uses a separate \emph{constraint register} connecting
Planck-scale configuration to field strength at locking. The
\emph{holographic width} at birth (\textasciitilde0.1 e-fold) is a third
coordinate --- horizon entropy at the transition, not the same as either
count above.

The universe counted its earliest ripples in one register and locked its
physics in another --- conflating the two was where the apparent tension
came from.

At birth itself, three clocks must not be mixed: thermal Friedmann scale
at \textasciitilde10³² K; slow constraint Hubble at locking; and
inflation-depth N\_e for the CMB. Same event, different projections.

\begin{center}\rule{0.5\linewidth}{0.5pt}\end{center}

\subsection{What Is Closed (and What Is
Not)}\label{what-is-closed-and-what-is-not}

Several foundations have tightened without withdrawing the headline
predictions:

\begin{itemize}
\tightlist
\item
  Cross-thickness gravity now follows from a single action (Janus
  action), reproducing κ = κ₀/(1 − κ₀).
\item
  Six chord slots read from one heat-kernel computation on the same
  geometric object.
\item
  Λ lands near 1.05× observed from the same e-fold depth as n\_s: loud
  vacuum term removed by the volume constraint, quiet residual
  renormalised at gauge birth. Caveat unchanged: α is partly anchored
  through present-day dark-energy data.
\item
  m\_\{ν₂\} is derived to \textasciitilde0.7\% of the solar mass scale
  with no observational input in the chain.
\item
  Two routes to n\_s are now one identity, read two ways.
\end{itemize}

Still open: the full SM mass hierarchy, first-generation neutrino
Yukawa, M\_R on the heat-kernel route, strong-field gravity near
horizons, and deriving α without observational anchoring.

\begin{center}\rule{0.5\linewidth}{0.5pt}\end{center}

\subsection{What We Do Not Yet Know}\label{what-we-do-not-yet-know}

\begin{itemize}
\tightlist
\item
  Specific quark and charged-lepton masses (ontology at gauge birth,
  numbers still inherited).
\item
  Full neutrino mass matrix beyond m\_\{ν₂\}.
\item
  Leptogenesis scale M₁ for η magnitude (existence and sign of asymmetry
  are geometric).
\item
  Residual lithium-7 factor (\textasciitilde1.58× after cross-thickness
  correction, not fully closed).
\item
  Strong-field gravity and horizons beyond static weak-field sign-off.
\end{itemize}

The list is shorter than a year ago; it's not empty.

\begin{center}\rule{0.5\linewidth}{0.5pt}\end{center}

\subsection{The Short Version}\label{the-short-version}

Photons experience zero proper time; massive particles do not. NTE reads
that correlation through a \textbf{boundary expansion field}
\textbar Φ\_NTE\textbar{} distinct from the SM Higgs, both set at
\textbf{gauge birth}. The universe is an expanding S³ boundary; \emph{c}
is expansion rate; we are embedded behind the photon wavefront.
Bilateral structure is two orientations of one thickened manifold.
Incomplete neutrino locking bridges orientations; \textbf{dark matter is
their cross-thickness gravity} with \textbf{κ ≈ 5.50}. Time's arrow is
geometric commitment at locking. Quantum probabilities are sphere
geometry after Hopf quotient. Octonions without time become quaternions
with time -\textgreater{} \textbf{SU(3)×SU(2)×U(1)}. Hubble tension and
DESI \emph{w} bias share void enhancement of \textbar Φ\_NTE\textbar².
Core predictions: no direct DM; Majorana neutrinos; three generations;
m\_\{ν₂\}, n\_s, θ₂₃ at structural precision; Λ and H₀ with stated
anchoring caveats.

\begin{center}\rule{0.5\linewidth}{0.5pt}\end{center}

\emph{Non-Temporal Emergence (NTE) --- the Null-Time Extension to the
Standard Model (G. White). \textbf{Popular Paper v10.8 v7 --- Reader
Edition.} Full Edition:
\texttt{Paper\_\textbackslash{}allowbreak\{\}Popular\_\textbackslash{}allowbreak\{\}v10\_\textbackslash{}allowbreak\{\}8\_\textbackslash{}allowbreak\{\}v6.md}.
Technical derivations: book Appendix B. Correspondence welcome.}

\begin{center}\rule{0.5\linewidth}{0.5pt}\end{center}

\end{document}
